\subsection*{The full relative pro-three Ferrari family}

\paragraph{FR.1 (Proved): the compact target and its congruence kernel.}
Retain the notation of FC.1--FC.3 and FP.40:
\[
 L=\widehat B_4/\overline{\langle z_4\rangle},\qquad
 C=\ker(L\longrightarrow\widehat B)
      =C_2\amalg C_2\amalg C_2.
\]
Let $O^3(C')$ be the intersection of the kernels of all continuous
homomorphisms from $C'$ to finite three-groups, and put
\[
 \Gamma=L/O^3(C'),\qquad
 C_\infty=C/O^3(C'),\qquad
 N=C'/O^3(C').
 \tag{FR.1a}
\]
The subgroup $C'$ is characteristic in $C$, which is normal in $L$;
$O^3(C')$ is characteristic in $C'$. Thus it is a closed normal
subgroup of $L$, and every group and quotient in (FR.1a) is a
genuine profinite group. There are exact sequences
\[
 1\longrightarrow N\longrightarrow C_\infty
 \longrightarrow Q=(C_2)^3\longrightarrow1,
 \qquad
 1\longrightarrow C_\infty\longrightarrow\Gamma
 \longrightarrow\widehat B\longrightarrow1.
 \tag{FR.1b}
\]
Here $N$ is free pro-three of rank five, since $C'$ is free
profinite of rank five. In particular $N$ is finitely generated,
and so are $C_\infty$ and $\Gamma$.

Write
\[
 M=N/\Phi(N)=N/(N'N^3)\simeq\mathbb F_3^5.
\]
The $Q$-action on $M$ has characters $110,101,011,111,111$
by FC.2. These characters span $Q^*$, so that action is faithful.
The centralizer of $N$ in $C_\infty$ is trivial: an element
centralizing $N$ acts trivially on $M$, hence has trivial image
in $Q$; it then belongs to the centre of the nonabelian free
pro-three group $N$, which is trivial. Here is a direct
proof of the free-pro-$p$ centre assertion used at every prime.
For a free pro-$p$ group of rank $r\geq2$, the universal map
\[
 x_i\longmapsto1+T_i
 \quad\hbox{into}\quad
 1+(T_1,\ldots,T_r)\mathbb F_p\langle\!\langle
 T_1,\ldots,T_r\rangle\!\rangle
\]
is injective. Indeed the target is pro-$p$ by truncation in
degree. For every finite $p$-group quotient $H$, substitution
$T_i\mapsto\bar x_i-1$ is defined in $\mathbb F_p[H]$
because its augmentation ideal is nilpotent. It recovers
the original quotient map, so an element in the displayed
kernel is trivial in every finite quotient. Nilpotence of
that augmentation ideal follows by induction on $|H|$:
choose a central element $h$ of order $p$; the central ideal
$(h-1)$ has $p$th power zero, and the quotient is the group
algebra of $H/\langle h\rangle$.

A formal series commuting with $T_1$ contains only powers
of $T_1$. To see this coefficientwise, moving trailing
$T_1$'s to the beginning preserves the coefficient by
$T_1A=AT_1$; any remaining word ending in another letter
has coefficient zero by the same identity. Thus the
centralizer of $T_1$ is $\mathbb F_p[[T_1]]$.
Its intersection with the centralizer of $T_2$ is
$\mathbb F_p$. A central group element has augmentation
one, so its series is one. This proves the asserted
centrelessness and hence
\[
           C_{C_\infty}(N)=1,
           \qquad Z(C_\infty)=1.
 \tag{FR.1c}
\]

Conjugation defines a continuous homomorphism
\[
 \rho:\Gamma\longrightarrow\operatorname{Aut}(C_\infty),
 \qquad \rho_0:\Gamma\longrightarrow\operatorname{GL}(M).
\]
The finite generation ensures that these automorphism groups
carry their usual profinite topology: finite characteristic
quotients are cofinal. All images below are compact, hence
closed. The subgroup $N$ is characteristic in $C_\infty$,
being its unique pro-three Sylow subgroup, so restriction gives
$\operatorname{Aut}(C_\infty)\to\operatorname{Aut}(N)$.
It is injective by (FR.1c): if an automorphism fixes $N$
pointwise, the difference between each element and its image
centralizes $N$.

The kernel of
$\operatorname{Aut}(N)\to\operatorname{GL}(M)$ is pro-three.
Here is a filtration proof. Set
\[
 N_1=N,\qquad N_{r+1}=\overline{N_r^3[N_r,N]}.
\]
These characteristic open subgroups are cofinal, because the
corresponding series terminates in every finite three-group.
The quotients $N_r/N_{r+1}$ are central elementary abelian
sections and are generated by the commutators and powers
coming from $N/N_2$. An automorphism trivial on $N/N_2$
therefore acts trivially on every such section. If $U_r$
denotes the subgroup of automorphisms trivial on $N/N_r$,
then, for $r\geq2$, the map
\[
 U_r/U_{r+1}\longrightarrow
 \operatorname{Hom}(N/N_2,N_r/N_{r+1}),\qquad
 \alpha\longmapsto\bigl(g\mapsto\alpha(g)g^{-1}\bigr)
 \tag{FR.1d}
\]
is an injective group homomorphism. Its target is elementary
abelian of exponent three. It follows by induction that
$U_2/U_r$ is a finite three-group for every $r$, and hence
$U_2=\varprojlim_rU_2/U_r$ is pro-three.
Consequently the congruence kernel
\[
 \ker\bigl(\rho(\Gamma)\longrightarrow\rho_0(\Gamma)\bigr)
                       \quad\hbox{is pro-three.}
 \tag{FR.1e}
\]
One may avoid the centreless restriction argument in proving
(FR.1e): the additional kernel is a group of cocycles from
$Q$ into $Z(N)$, itself pro-three. The centreless statement
is needed later for the stronger conclusion inside $\Gamma$.

\paragraph{FR.2 (Proved): the exact braid centre controls both harmonic pairs.}
Let $G$ be the closed subgroup of $\operatorname{Aut}(C_\infty)$
generated by $\rho(a),\rho(b)$, and let $G_0$ be its image
on $M=B\oplus P$. On the boundary space $B$, its image
is $S_3$; on the proper space $P$, it is
$\operatorname{SL}_2(\mathbb F_3)$, with
\[
 a\longmapsto U=\begin{pmatrix}1&1\\0&1\end{pmatrix},
 \qquad
 b\longmapsto L_0=\begin{pmatrix}1&0\\-1&1\end{pmatrix}.
 \tag{FR.2a}
\]
The boundary assertion follows either from the exact actions
(FC.2a), or the cube certificate: the two boundary generators
are involutions satisfying the braid relation and generate
a group of order six. Thus $G_0$ embeds in
$S_3\times\operatorname{SL}_2(\mathbb F_3)$ and projects
onto both factors. Since
\[
 S_3''=1,\qquad
 \operatorname{SL}_2(\mathbb F_3)'=Q_8,\qquad
 \operatorname{SL}_2(\mathbb F_3)''=\{\pm I\},
\]
its second derived group is exactly
\[
              G_0''=\{(I_B,I_P),(I_B,-I_P)\}\simeq C_2.
 \tag{FR.2b}
\]
No assertion about independence of the two linear factors is
needed for this conclusion: the inclusion follows from the
two displayed derived groups, and surjectivity onto the proper
factor proves equality.

The source element $z_3=(ab)^3$ is central in the actual
three-strand section. Its image $z=\rho(z_3)$ is therefore
central in $G$. On $M$ it has the value $(I_B,-I_P)$.
Let $e_2\in\widehat{\mathbb Z}$ be one in the two-adic
factor and zero in all other prime factors. Since the kernel
of $G\to G_0$ is pro-three, the procyclic group generated
by $z$ has order dividing $2\cdot3^\infty$ and maps onto
the displayed $C_2$. Hence
\[
                         \iota=z^{e_2}
\]
is a central involution of $G$ with the same nontrivial
linear image. If $U_G=\ker(G\to G_0)$, then
\[
  (G\longrightarrow G_0)^{-1}(G_0'')
                 =U_G\times\langle\iota\rangle.
 \tag{FR.2c}
\]
The right-hand side is pronilpotent, being the direct product
of a pro-three group and a group of order two. It contains
$G''$. Every closed subgroup of a pronilpotent group is
pronilpotent. Thus the second derived group of
$\overline{\langle\rho(a^2),\rho(b^2)\rangle}$ is
pronilpotent. Reversing the braid generators gives the identical
conclusion for
$\overline{\langle\rho(b^2),\rho(c^2)\rangle}$.

This central lift is essential. An arbitrary extension of
$C_2$ by a pro-three group need not be pronilpotent. The
argument uses centrality of the literal braid $z_3$ in
the actual harmonic subgroup; it does not assume that
$-I$ acts trivially on every direction in the ambient
automorphism group.

\paragraph{FR.3 (Proved): all five pair second-derived images are pronilpotent.}
Use the five literal pentagon inputs of FP.36 and FP.41:
\[
 X_1=a^2,\quad X_2=b^2,\quad X_3=c^2,\quad
 X_4=(aba)^2=(ab)^3,\quad X_5=(cbc)^2=(cb)^3.
 \tag{FR.3a}
\]
Every $X_i$ has identity boundary linear action. On the
proper two-dimensional block their images are respectively
\[
                    U^2,\quad L_0^2,\quad U^2,\quad -I,\quad -I.
 \tag{FR.3b}
\]
The two harmonic pairs $(X_1,X_2)$ and $(X_2,X_3)$
are covered by FR.2. Each of the other three pairs
$(X_3,X_4)$, $(X_5,X_1)$, $(X_4,X_5)$ has abelian
linear image, by (FR.3b). Its first derived action is
therefore contained in the pro-three congruence kernel
of (FR.1e). In particular its second derived action
is pro-three. All five action-pair groups consequently
have pronilpotent second derived subgroup.

Let $R_{\rm nil}(F'')$ be the kernel of the maximal
pronilpotent quotient of the second derived subgroup of
the free profinite group $F=\widehat F(x,y)$. For every
\[
                         f\in R_{\rm nil}(F''),
 \tag{FR.3c}
\]
the image of $f$ under each of the five pair evaluations
in $\operatorname{Aut}(C_\infty)$ is one. Indeed a pair
map sends $F''$ onto the closed second derived group
of its image. A continuous map from $F''$ to a
pronilpotent group kills $R_{\rm nil}(F'')$, by its
defining universal property.

This conclusion is simultaneous at every finite characteristic
three-group quotient of $C'$ and in their inverse limit.
No selection of finite residues or passage from finite
solutions to a genuine $\widehat{GT}$ element occurs.
Adjoining the forgetful $q_{12}$ action does not change it:
the five input images there are $(x,y,x,1,1)$; the two
nontrivial pair groups have nilpotent second derived
group. Thus the joint pair second-derived image is a
closed subgroup of a product of pronilpotent groups.

\paragraph{FR.4 (Proved, with LI.8): the pentagon holds in the actual quotient with full forgetful base.}
Suppose in addition to (FR.3c) that
$f(y,x)=f(x,y)^{-1}$. Let $\mathcal P(f)$ denote the
literal five-point pentagon product in $L$. Its image
under $L\to\widehat B$ is one: the five inputs become
$(x,y,x,1,1)$, so the only nontrivial factors cancel
by this two-cycle identity, as in FP.30. Consequently
the image of $\mathcal P(f)$ in $\Gamma=L/O^3(C')$
belongs to $C_\infty$.

By FR.3, each factor has trivial conjugation action on
$C_\infty$. Their product therefore centralizes
$C_\infty$. Equation (FR.1c) now gives
\[
             \boxed{\ \mathcal P(f)=1
                       \quad\hbox{in }L/O^3(C').\ }
 \tag{FR.4a}
\]
This is stronger than equality in an automorphism image:
the quotient retains the complete profinite forgetful
base $\widehat B$. Individual factors are asserted to
have trivial action; they need not themselves be one
in this larger quotient. It is their pentagon product
that belongs to the centreless kernel and hence is one.

LI.8 supplies an actual gentle element $k$ with
\[
            f_k\in R_{\rm nil}(F''),\qquad\Delta(k)=1.
\]
Its two-cycle identity is part of the genuine gentle
equations. Thus this same actual element satisfies
(FR.4a) and all of its finite relative three-group
quotients, while retaining $\Delta(k)=1$. The entire
relative pro-three Ferrari family is therefore
insufficient to force the required parity identity.
This does not assert the full profinite pentagon,
membership of $k$ in $\widehat{GT}$, or an analogous
result for relative primes other than three.

\paragraph{FR.5 (Proved): the integral marking and the other relative primes.}
For every prime $\ell$, put
\[
 C_\ell=C/O^\ell(C'),\qquad
 \Gamma_\ell=L/O^\ell(C').
 \tag{FR.5a}
\]
For odd $\ell$, the cube idempotents of FC.2 are defined over
$\mathbb Z[1/2]$. More intrinsically, averaging over the deck
group $K/C'$ of order four splits the integral homology after
inverting two into puncture homology and proper elliptic homology.
The proper block is identified with $H_1(K)$ by pushforward
and transfer, whose composite is multiplication by four.
The actions of $a,b,c$ on that block are respectively $U,L_0,U$
over $\mathbb Z$. On the boundary block each pure $X_i$ acts
trivially: $a,b,c$ permute the four two-torsion punctures,
their squares are trivial on them, and each of $(ab)^3,(cb)^3$
has trivial permutation. Thus for every odd $\ell$ the five
mod-$\ell$ linear actions are exactly
\[
        I_B\oplus U^2,\quad I_B\oplus L_0^2,\quad
        I_B\oplus U^2,\quad I_B\oplus(-I),\quad
        I_B\oplus(-I).
 \tag{FR.5b}
\]
This assertion is not obtained by extrapolating from characteristic
three. The integral part of \texttt{check\_ferrari\_cube.py}
constructs the integer homology matrices and verifies
$X_i(4I-\sum_{q\in K/C'}q)=4I-\sum_{q\in K/C'}q$ for
each $i$, together with the five exact pushforward identities
on the proper homology. These are integer equalities, and
division by four is legitimate at every odd prime.

The same proof as FR.1, with three replaced by $\ell$, shows
that the kernel of the action on $H_1(C',\mathbb F_\ell)$
is pro-$\ell$. Consequently the procyclic action groups of
$X_1,X_2,X_3$ are pro-$\ell$, and those of $X_4,X_5$
are pro-$\{2,\ell\}$. This is only a statement about each
individual input; no assertion that their generated pair
group is pro-$\{2,\ell\}$ is needed.

At $\ell=2$ there is the exact identification
\[
                   C_2=C/O^2(C')=C^{(2)}.
 \tag{FR.5c}
\]
Indeed the left-hand side is an extension of a pro-two group
by $(C_2)^3$, hence is pro-two; every finite two-group quotient
of $C$ kills $O^2(C')$, proving its universal property.
Its abelianization is $(C_2)^3$, and its Frattini quotient is
the same group, since $C$ is generated by involutions. The
exact actions (FC.2a) show that each pure $X_i$ acts trivially
on this quotient. The Frattini automorphism kernel is pro-two
by the filtration argument of FR.1, which applies to every
finitely generated pro-$p$ group. Thus all five input actions
are pro-two in this case.

Each $C_\ell$ is centreless. For odd $\ell$ this follows
as in FR.1 from the faithful $(C_2)^3$ action on the rank-five
mod-$\ell$ homology and the centrelessness of the free pro-$\ell$
kernel. For two, the split presentation $C=K\rtimes\langle j\rangle$
gives
\[
                  C^{(2)}=K^{(2)}\rtimes\langle j\rangle,
\]
where $K^{(2)}$ is free pro-two of rank two. An element of
$K^{(2)}$ central in this extension is one, and an element
outside it induces $-I$ on its abelianization $\mathbb Z_2^2$,
so cannot be central. This proves the required centreless
assertion without reducing $-I$ modulo two.

\paragraph{FR.6 (Proved, using R2): one actual witness for the full relative pronilpotent quotient.}
Let $e\in\widehat{\mathbb Z}$ have $e_3=1$ and $e_p=0$
for every $p\ne3$, and set
\[
                 A_e=[[x^e,y^e],[y^e,z^e]],
                 \qquad z=(xy)^{-1}.
 \tag{FR.6a}
\]
For every $\ell\ne3$, (FR.5) implies that each of the five
$X_i^e$ is one in $\operatorname{Aut}(C_\ell)$. Thus $A_e$
is in the kernel of every pentagon pair evaluation at that
prime. It is in each required peripheral conjugate of these
kernels as well: the peripheral $S_3$ permutes, up to inner
conjugacy, the three source peripheral elements $x,y,z$;
the original two input powers are one, and if any one of the
three powered entries in (FR.6a) is one, its double commutator
is one. Passing to the finite $\widehat B$-normal core therefore
preserves membership of $A_e$.

The same seed lies in $F''$, has $e\equiv4\pmod{12}$ and
$q_{12}(A_e)=5I$, and lies in the full original orbifold cores:
one source peripheral has order two there and $e$ is even.
These are precisely the seed conditions used in LI.8--LI.9.
For any finite list of finite action quotients at primes
$\ell\ne3$, intersect their peripheral normal cores with
the finite normal subgroups used to impose the LI.8 conditions.
The single seed $A_e$ remains in this intersection and retains
its prescribed signed value. The relative R2 construction of
LI.9 consequently supplies a genuine gentle element satisfying
that finite list, with $\Delta=1$ and the corresponding
relative pronilpotent conditions.

All these conditions are closed subsets of the compact group
of genuine gentle elements: finite evaluation equations are
clopen, $\Delta=1$ is clopen, and the relative pronilpotent
conditions are closed inverse-image conditions. The finite
intersection property, with all open normal subgroups needed
in LI.8 included in the indexing family, therefore gives one
actual element $k_{\rm rel}$ such that
\[
 \Delta(k_{\rm rel})=1,\qquad
 f_{k_{\rm rel}}\in R_{\rm nil}(F''),
 \qquad
 f_{k_{\rm rel}}(X_i,X_j)=1
       \text{ in }\operatorname{Aut}(C_\ell)
       \quad(\ell\ne3)
 \tag{FR.6b}
\]
for every one of the five pairs. The prime-three assertion
then follows from FR.3, using the same actual element.
This does not assert that the particular earlier LI.8 witness
already satisfies all the new prime conditions; it constructs
a simultaneous strengthening by the same authorized R2 theorem.

The entire LI.9 Eisenstein--Burau family may be imposed in
the same finite-intersection construction. At any finite list
of Burau congruence levels, the exponent $e$ reduces to an
integer $m$ divisible by every required prime-to-three
exponent and congruent to one modulo a sufficiently large
power of three. In particular $m$ is even, prime to three,
and congruent to four modulo twelve. These are the allowed
LI.9 seed exponents. Its square-zero peripheral identities
therefore kill the same nonharmonic pair kernels away from
three; its harmonic and three-adic conditions follow from
$f\in R_{\rm nil}(F'')$ exactly as in LI.9. Thus the fixed
seed $A_e$ retains all LI.9 conditions simultaneously with
the new Ferrari conditions, rather than requiring two
unrelated witnesses.

For each prime, antisymmetry puts the pentagon defect inside
$C_\ell$, while (FR.6b) and FR.3 make its inner action trivial.
The centreless result in FR.5 proves
$\mathcal P(f_{k_{\rm rel}})=1$ in every $\Gamma_\ell$.
Finally
\[
       \bigcap_\ell O^\ell(C')=R_{\rm nil}(C'),
 \tag{FR.6c}
\]
because a finite nilpotent group is the direct product of its
Sylow subgroups. The natural map
$L/R_{\rm nil}(C')\to\prod_\ell L/O^\ell(C')$ is therefore
injective, as well as continuous with compact closed image.
Consequently
\[
 \boxed{\quad\mathcal P(f_{k_{\rm rel}})=1
           \text{ in }L/R_{\rm nil}(C'),\qquad
           \Delta(k_{\rm rel})=1.\quad}
 \tag{FR.6d}
\]
The quotient retains the full profinite forgetful base and
the full relative pronilpotent Ferrari kernel. This proves
the insufficiency of this entire family, conditional only
on the already specified R2 lifting input used in LI.8--LI.9.
It is not the full five-point pentagon, and it does not
produce an element of $\widehat{GT}$ with nonzero discrepancy.

\paragraph{FR.7 (Proved, with the same R2 witness): an actual automorphism of the relative quotient.}
There is a further action-level conclusion that does not assume
kernel preservation merely from the pentagon product. Put
\[
 \Gamma_{\rm nil}=L/R_{\rm nil}(C'),\qquad
 C_{\rm nil}=C/R_{\rm nil}(C').
\]
The group $C_{\rm nil}$ is centreless. Indeed it embeds into
$\prod_\ell C_\ell$, and its projection to each $C_\ell$
is surjective. A central element therefore has trivial image
in every centreless $C_\ell$, and is one.

For the witness $k_{\rm rel}=(1,f)$ of FR.6, define in
$\Gamma_{\rm nil}$
\[
 f_0=f(a^2,b^2)=f(X_1,X_2),\qquad
 e=f(z_3,c^2)=f(X_4,X_3),
\]
and assign
\[
       a\longmapsto a'=a,\qquad
       b\longmapsto b'=f_0^{-1}bf_0,\qquad
       c\longmapsto c'=e^{-1}ce.
 \tag{FR.7a}
\]
These are exactly the multiplier-one braid formulas in
the marking of FP.26.4. By FR.6, $f_0$ acts trivially on
every $C_\ell$; the same is true of $e$, using antisymmetry
to reverse the pair $(X_3,X_4)$. They therefore centralize
$C_{\rm nil}$. Hence the three new generators have exactly
the same conjugation action on $C_{\rm nil}$ as the old
generators.

Under the fold $q:\Gamma_{\rm nil}\to\widehat B$, one has
$q(c)=q(a)$ and $q(z_3)=1$, so $q(e)=f(1,q(c)^2)=1$.
Consequently their folded images are
\[
 q(a')=\Phi_3(q(a)),\qquad
 q(b')=\Phi_3(q(b)),\qquad
 q(c')=\Phi_3(q(a)),
 \tag{FR.7b}
\]
where $\Phi_3$ is the existing genuine gentle automorphism
of the three-strand central quotient in this marking.
That source automorphism, not any five-point assumption,
is the input to (FR.7b).

Now evaluate any of the four defining relators
\[
 aba=bab,\qquad bcb=cbc,\qquad ac=ca,\qquad(abc)^4=1
\]
on $(a',b',c')$. Its fold is one by (FR.7b), and its
conjugation action on $C_{\rm nil}$ is one because the
three actions have not changed. The relator therefore
belongs to $Z(C_{\rm nil})=1$. The assignment (FR.7a)
thus defines a continuous homomorphism
$\widetilde\Phi:L\to\Gamma_{\rm nil}$ directly from
the profinite presentation.

For every $w\in C$, both $\widetilde\Phi(w)$ and its
original quotient image $\bar w$ belong to $C_{\rm nil}$,
and they induce the same inner automorphism there.
Centrelessness gives $\widetilde\Phi(w)=\bar w$.
In particular $R_{\rm nil}(C')$ is killed. Thus the map
descends to an endomorphism
\[
            \Phi_{\rm nil}:\Gamma_{\rm nil}
                         \longrightarrow\Gamma_{\rm nil}
\]
that fixes $C_{\rm nil}$ pointwise and induces the
automorphism $\Phi_3$ on the quotient $\widehat B$.
It is injective: an element of its kernel first maps
trivially to $\widehat B$, hence lies in the fixed kernel
$C_{\rm nil}$. It is surjective: lift a preimage under
$\Phi_3$ of the desired folded element, and then correct
the difference by an element of the fixed kernel.
A continuous bijection of compact Hausdorff groups is
a topological isomorphism. Therefore
\[
 \boxed{\ \Phi_{\rm nil}\in\operatorname{Aut}(\Gamma_{\rm nil}),
 \qquad\Phi_{\rm nil}|_{C_{\rm nil}}=\operatorname{id},
 \qquad\Delta(k_{\rm rel})=1.\ }
 \tag{FR.7c}
\]
This proves kernel stability and an actual relative action.
It does not promote $k_{\rm rel}$ to an element of the
full profinite Grothendieck--Teichm\"uller group.
